\documentclass[a4paper]{article}

\usepackage[toc,page]{appendix}

\usepackage[spanish]{babel}
\usepackage[utf8]{inputenc}
\usepackage{amsmath}
\usepackage{graphicx}
\usepackage{listings}
\usepackage{fancyhdr}
\usepackage{amsmath}
\usepackage[colorinlistoftodos]{todonotes}
\usepackage{xcolor}
\usepackage{minted}
\usepackage{enumitem}
\usepackage{hyperref}

\hypersetup{
    colorlinks,
    citecolor=blue,
    filecolor=blue,
    linkcolor=blue,
    urlcolor=blue
}

\definecolor{mGreen}{rgb}{0,0.6,0}
\definecolor{mGray}{rgb}{0.5,0.5,0.5}
\definecolor{mPurple}{rgb}{0.58,0,0.82}
\definecolor{backgroundColour}{rgb}{0.95,0.95,0.92}



\usepackage{geometry}
 \geometry{
 a4paper,
 total={210mm,297mm},
 left=22mm,
 right=19.1mm,
 top=25.4mm,
 bottom=25.4mm,
 }
 
 \lstdefinestyle{CStyle}{
    backgroundcolor=\color{backgroundColour},   
    commentstyle=\color{mGreen},
    keywordstyle=\color{magenta},
    numberstyle=\tiny\color{mGray},
    stringstyle=\color{mPurple},
    basicstyle=\footnotesize,
    breakatwhitespace=false,         
    breaklines=true,                 
    captionpos=b,                    
    keepspaces=true,                 
    numbers=left,                    
    numbersep=5pt,                  
    showspaces=false,                
    showstringspaces=false,
    showtabs=false,                  
    tabsize=2,
    language=C
}


\begin{document}
\begin{figure}
\centering
\includegraphics[scale=1]{./img/logo-facu}
\end{figure}  

\title{\large\textsc{66.20 - Organización de Computadoras}\\
\large Trabajo Práctico 2 - Memorias Caché}

\author{
Andrew Parlane\\
}

\maketitle

\newpage

\tableofcontents

\newpage

\section{Resumen}
El objetivo de este trabajo se basa en aplicar los conocimientos aprendidos en clase acerca caches. Para eso desarrollé un programa en C que simula un caché.

\section{Introducción}
El programa que desarrollé lee comandos desde un archivo texto y simula lecturas y escrituras a memoria usando un cache \textit{Write Through / No Write Allocate}. Hay tres comandos soportado \textit{R} para \textit{Read}, \textit{W} para \textit{Write} y \textit{MR} para \textit{Miss Rate}. La memoria principal es 4KB, y el cache es 1KB asociativa por conjuntos con dos vías. El tamaño de bloque es 32 bytes.

\section{Compilacion y ejecución del programa}
Esta sección describe  los mecanismos para la compilación y ejecución del programa.

Para compilar usando el Makefile:
\begin{itemize}
\item Estando parados en la carpeta donde se encuentran los archivos fuente, ejecutamos el siguiente comando:
\begin{minted}{bash}
make
\end{minted}
\end{itemize}

Para ejecutar:
\begin{itemize}
\item Estando parados en la carpeta donde se encuentran el archivo ejecutable, corremos para cada archivo lo siguiente:
\begin{minted}{bash}
./tp2 [options] file 
\end{minted}
\item Si se quiere ver el help, que especifica las opciones disponibles al momento de invocar el programa:
\begin{minted}{bash}
./tp2 -h
\end{minted}
\item Si se quiere ver la versión del programa:
\begin{minted}{bash}
./tp2 -V
\end{minted}
\end{itemize}

\section{Implementación}
\subsection{Implementación en C}

Usé algunas definiciones y estructuras para simplificar la implementación de memoria principal y el caché:

\begin{itemize}
\item \textit{Bloque} - El bloque es un struct que contiene un array de TAMANO\_DE\_BLOQUE bytes.
\item \textit{LineaDeCache} - Una línea de caché contiene un Bloque de data, el TAG y un flag válido.
\item \textit{ViaDeCache} - Una vía de caché contiene CONJUNTOS\_EN\_CACHE líneas.
\item \textit{Cache} - El cache contiene NUMERO\_DE\_VIAS vías y un array de CONJUNTOS\_EN\_CACHE de tipo bool que uso por el algoritmo LRU.
\end{itemize}

A continuación enumero las funciones definidas en el programa que se usarán luego de la verificación y validación de los parámetros de entrada: 
\begin{itemize}
\item \textbf{init}\\
	\begin{tabular}{l p{12cm}}
    parámetros: &\\
    &\\
    descripción: &función que inicializa el cache, escribiendo cero al bit válido de todos las líneas del cache.
    \end{tabular}\\

\item \textbf{dumpCache}\\
	\begin{tabular}{l p{12cm}}
    parámetros: &\\
    &\\
    descripción: &función que escribe el contenido del cache a stdout. Solo usado si DEBUG es definido.
    \end{tabular}\\

\item \textbf{decodeAddress}\\
	\begin{tabular}{l p{12cm}}
    parámetros: &uint32\_t addr \\
                &uint32\_t *tag \\
                &uint32\_t *indice \\
                &uint32\_t *via \\
                &uint32\_t *offset \\
                &uint32\_t *mpBloque \\
    &\\
    descripción: &función que convierta una dirección de memoria principal hasta el índice de conjunto, el TAG, el número de bloque en MP y el offset del byte en ese bloque. Después busca en los dos vías del cache a ver si el byte querido es en el caché. Si no es encontrado incrementa el contador de misses.
    \end{tabular}\\

\item \textbf{read\_byte}\\
	\begin{tabular}{l p{12cm}}
    parámetros: &int address \\
    &\\
    descripción: &función que lee el byte a la dirección: \textit{address}. Si no es en el caché copia el bloque querido desde MP hasta el caché, reemplazando la línea usado más lejos en el pasado.
    \end{tabular}\\

\item \textbf{write\_byte}\\
	\begin{tabular}{l p{12cm}}
    parámetros: &int address \\
    			&int value \\
    &\\
    descripción: &función que escribe el byte: \textit{value} a la dirección: \textit{address} en MP. Si es en caché actualiza el valor en el cache también.
    \end{tabular}\\
    
\item \textbf{get\_miss\_rate}\\
	\begin{tabular}{l p{12cm}}
    parámetros: &\\
    &\\
    descripción: &función que devuelva el \textit{Miss Rate} como un porcentaje redondeado al entero más próximo. Si no había accesos al caché devuelva 0\%.
    \end{tabular}\\
    
\item \textbf{usage}\\
	\begin{tabular}{l p{12cm}}
    parámetros: &FILE *stream \\
    			&const char *nuestroNombre\\
    &\\
    descripción: &función que muestra el help de la aplicación.
    \end{tabular}\\

\item \textbf{skipLine}\\
	\begin{tabular}{l p{12cm}}
    parámetros: &FILE *stream \\
    &\\
    descripción: &función que bota todos los caracteres en el archivo hasta el final de la línea corriente. Uso esto cuando había un error en un comando.
    \end{tabular}\\

\item \textbf{leerLongLong}\\
	\begin{tabular}{l p{12cm}}
    parámetros:    &FILE *f\\
                   &long long *ll\\
                   &bool *OK\\
                   &bool *eof\\
                   &bool *newLine\\
	&\\
    descripción: &función que lee un entero de desde el archivo. Lo copié desde TP1.
    \end{tabular}\\
    
\item \textbf{leerComando}\\
	\begin{tabular}{l p{12cm}}
    parámetros:    &FILE *f\\
                   &bool *eof\\
                   &Comando *cmd\\
	&\\
    descripción: &función que lee el comandos siguiente desde el archivo.
    \end{tabular}\\

\item \textbf{leerArchivo}\\
	\begin{tabular}{l p{12cm}}
    parámetros:    &const char *archivo \\
	&\\
    descripción: &función que lee el archivo, y actúa sobre los comandos encontrados.
    \end{tabular}\\

\end{itemize}
\lstset{basicstyle=\footnotesize\ttfamily}

\section{Pruebas}
Realicé las pruebas incluido con el enunciado. 
\setlist{nosep}
\begin{itemize}
\item \textbf{prueba1.mem}
\item \textbf{prueba2.mem}
\item \textbf{prueba3.mem}
\end{itemize}

\begin{minted}{bash}
$ ./tp2 pruebas/prueba1.mem
Miss Rate: 67%

$ ./tp2 pruebas/prueba2.mem
Miss Rate: 50%

$ ./tp2 pruebas/prueba3.mem
Miss Rate: 60%

$ ./tp2 pruebas/prueba4.mem
Miss Rate: 53%

$ ./tp2 pruebas/prueba5.mem
Miss Rate: 100%
\end{minted}

También He realizado unas pruebas extras con archivos mal formados.

\section{Conclusiones}
	El desarrollo de este trabajo me permitió llevar a la práctica los conocimientos adquiridos acerca de un caché asociativa por conjuntos WT/nWA.

\section{Código}

\begin{minted}{C}
#include <stdio.h>
#include <stdlib.h>
#include <libgen.h>
#include <getopt.h>
#include <ctype.h>
#include <stdint.h>
#include <stdbool.h>
#include <math.h>
#include <errno.h>
#include <string.h>

#define MAJOR_VERSION   1
#define MINOR_VERSION   1

//#define DEBUG

// No pueden cambiar estos sin cambiando el código
#define TAMANO_DE_BLOQUE    (32)
#define TAMANO_DE_MEMORIA   (4096)
#define TAMANO_DE_CACHE     (1024)
#define NUMERO_DE_VIAS      (2)

// Valores calculados
#define BLOQUES_EN_MP       ((TAMANO_DE_MEMORIA)/(TAMANO_DE_BLOQUE))
#define LINEAS_EN_CACHE     ((TAMANO_DE_CACHE)/(TAMANO_DE_BLOQUE))
#define CONJUNTOS_EN_CACHE  ((LINEAS_EN_CACHE)/(NUMERO_DE_VIAS))

typedef enum
{
    Comando_R = 0,
    Comando_W,
    Comando_MR,
} Comando;

typedef struct
{
    uint8_t data[TAMANO_DE_BLOQUE];
} Bloque;

typedef struct
{
    Bloque data;

    uint32_t tag;
    bool valido;
} LineaDeCache;

typedef struct
{
    LineaDeCache linea[CONJUNTOS_EN_CACHE];
} ViaDeCache;

typedef struct
{
    ViaDeCache via[NUMERO_DE_VIAS];

    bool via1estuvoUsadoUltimo[CONJUNTOS_EN_CACHE];
} Cache;

static const struct option _gLongOptions[] =
{
    {"help",    no_argument,        0, 'h' },
    {"version", no_argument,        0, 'V' },
    {0,         0,                  0,  0  }
};

static Bloque   _gMemoriaPrincipal[BLOQUES_EN_MP];
static Cache    _gCache;

static uint32_t _gAccessosTotal;
static uint32_t _gMisses;

// ------------------------------------------------------------------
// Funciones obligatorio por el TP
// no cambias los nombres, argumentos
// o tipo de resultados
// ------------------------------------------------------------------

static void init()
{
    _gAccessosTotal = 0;
    _gMisses = 0;

    uint32_t via;
    for (via = 0; via < NUMERO_DE_VIAS; via++)
    {
        uint32_t linea;
        for (linea = 0; linea < CONJUNTOS_EN_CACHE; linea++)
        {
            _gCache.via[via].linea[linea].valido = false;
        }
    }
}

#ifdef DEBUG
static void dumpCache()
{
    printf("IDX    VIA1    VIA2    Ultimo\n");
    int i;
    for (i = 0; i < (CONJUNTOS_EN_CACHE); i++)
    {
        printf("%3X    ", i);
        if (_gCache.via[0].linea[i].valido)
        {
            printf("%4X    ", _gCache.via[0].linea[i].tag);
        }
        else
        {
            printf("INVA    ");
        }
        if (_gCache.via[1].linea[i].valido)
        {
            printf("%4X    ", _gCache.via[1].linea[i].tag);
        }
        else
        {
            printf("INVA    ");
        }
        printf("%u\n", _gCache.via1estuvoUsadoUltimo[i]);
    }
    printf("\n");
}

#endif

static bool decodeAddress(uint32_t addr, uint32_t *tag, uint32_t *indice, uint32_t *via, uint32_t *offset, uint32_t *mpBloque)
{
    // Hay 32 bytes cada bloque, así offset = 5 bits
    // Hay 16 conjuntos, así indice = 4 bits
    // Y TAG es los demás
    // |---MPBLOQUE----|--OFFSET--| <-- MP
    // |---TAG---|-IDX-|--OFFSET--| <-- Cache
    //  15      9 8   5 4        0

    *offset = addr & 0x1F;
    *indice = (addr >> 5) & 0x0F;
    *tag = (addr >> 9);
    *mpBloque = (addr >> 5);

#ifdef DEBUG
    printf("addr %X:\n", addr);
    printf("  tag: %X\n", *tag);
    printf("  idx: %X\n", *indice);
    printf("  off: %X\n", *offset);
    printf("  mp:  %X\n", *mpBloque);
#endif

    // Comproba si es un HIT
    LineaDeCache *linea0 = &_gCache.via[0].linea[*indice];
    LineaDeCache *linea1 = &_gCache.via[1].linea[*indice];

    if (linea0->valido &&
        linea0->tag == *tag)
    {
        // HIT en via 0
        *via = 0;
#ifdef DEBUG
        printf("  HIT\n\n");
#endif
        return true;
    }

    if (linea1->valido &&
        linea1->tag == *tag)
    {
        // HIT en via 0
        *via = 1;
#ifdef DEBUG
        printf("  HIT\n\n");
#endif
        return true;
    }

    // MISS
#ifdef DEBUG
    printf("  MISS\n\n");
#endif
    _gMisses++;
    return false;
}

static int read_byte(int address)
{
    _gAccessosTotal++;

    uint32_t indice;
    uint32_t tag;
    uint32_t via;
    uint32_t offset;
    uint32_t mpBloque;

    unsigned char data;
    bool hit = decodeAddress(address, &tag, &indice, &via, &offset, &mpBloque);
    if (hit)
    {
        // hit
    }
    else
    {
        // miss

        // si hay un linea que no es válido reemplazamos eso
        if (!_gCache.via[0].linea[indice].valido)
        {
            via = 0;
        }
        else if (!_gCache.via[1].linea[indice].valido)
        {
            via = 1;
        }
        else
        {
            // LRU
            via = !_gCache.via1estuvoUsadoUltimo[indice];

            // Cache es WT así no puede ser dirty
        }

        memcpy(_gCache.via[via].linea[indice].data.data,
               _gMemoriaPrincipal[mpBloque].data,
               TAMANO_DE_BLOQUE);

        _gCache.via[via].linea[indice].tag = tag;
        _gCache.via[via].linea[indice].valido = true;
    }

    // leer Data de cache
    data = _gCache.via[via].linea[indice].data.data[offset];

#ifdef DEBUG
    printf("Read byte %u\n", data);
#endif

    // actualiza ultimo uso
    _gCache.via1estuvoUsadoUltimo[indice] = via;

    return hit ? data : -1;
}

static int write_byte(int address, unsigned char value)
{
    _gAccessosTotal++;
    uint32_t indice;
    uint32_t tag;
    uint32_t via;
    uint32_t offset;
    uint32_t mpBloque;

    bool hit = decodeAddress(address, &tag, &indice, &via, &offset, &mpBloque);
    if (hit)
    {
        // hit
        // actualizar el cache
        _gCache.via[via].linea[indice].data.data[offset] = value;

        // actualiza ultimo uso
        _gCache.via1estuvoUsadoUltimo[indice] = via;
    }
    else
    {
        // miss
        // ~WA así no hacemos nada
    }

    // actualizar MP (estamos WT)
    _gMemoriaPrincipal[mpBloque].data[offset] = value;

    return hit ? 0 : -1;
}

static unsigned int get_miss_rate()
{
    if (_gAccessosTotal == 0)
    {
        // _gMisses / 0
        // _gMisses debería estar 0 también
        // así decimos 0%
        return 0;
    }
    else
    {
        double porcentaje = ((double)_gMisses * 100.0l) / (double)_gAccessosTotal;
        return (unsigned int)round(porcentaje);
    }
}

// ------------------------------------------------------------------
// Final de funciones obligatorio por el TP
// ------------------------------------------------------------------

static void usage(FILE *stream, const char *nuestroNombre)
{
    fprintf(stream,
           "Usage:\n"
           "  %s -h\n"
           "  %s -V\n"
           "  %s filename\n"
           "Options:\n"
           "  -h, --help Prints usage information.\n"
           "  -V, --version Prints version information.\n"
           "\n\n",
           nuestroNombre, nuestroNombre, nuestroNombre);
    // necesitamos usar nuestroNombre, nuestroNombre, nuestroNombre
    // porque no soportamos %1$s
}

static void skipLine(FILE *f)
{
    while (1)
    {
        int res = fgetc(f);
        if (res == EOF ||
            (char)res == '\n')
        {
            return;
        }
    }
}

// lea carácter por carácter deshaciendo whitespace
// hasta encontrar [0-9-]. Después comenzar a leer números
// [0-9]. hasta obtenemos EOF, \n, \r, ' ', \t, ','
// Es un error si encontramos algún otro carácter.
// devolver no 0 si hay un error
// *OK = 1 -> hay un integer valido en *ll
// *eof = 1 -> no hay más a leer
// *newLine = 1 -> encontramos nueva línea
static bool leerLongLong(FILE *f, long long *ll, bool *OK, bool *eof, bool *newLine)
{
    *OK = false;
    *eof = false;
    *newLine = false;

    // soportamos signed 64 bits:
    // máx = 0x7FFF_FFFF_FFFF_FFFF =  9223372036854775807
    // min = 0x8000_0000_0000_0000 = -9223372036854775808
    // así max input legal es 20 cáracters +1 por NULL terminator
#define MAX_CHARS 20
    char buff[MAX_CHARS + 1];
    uint idx = 0;

    bool comenzandoLeerInt = false;
    while (1)
    {
        int res = fgetc(f);
        if (res == EOF)
        {
            *eof = true;
            // si tenemos algo en buff, convertimos ahora
            if (*OK)
            {
                buff[idx] = '\0';
                *ll = strtoll(buff, NULL, 10);
                if (errno != 0)
                {
                    fprintf(stderr, "Failed to convert %s to long long, error: %s\n", buff, strerror(errno));
                    return false;
                }
                return true;
            }
            else if (comenzandoLeerInt)
            {
                // solo podríamos estar aquí si leemos
                // un '-' y después nada, eso es un error
                fprintf(stderr, "Found invalid entry \"-\"\n");
                return false;
            }
            else
            {
                // eof but no error
                return true;
            }
        }

        char c = (char)res;
        if (c == '\r' || c == '\n')
        {
            *newLine = true;
        }

        if (!comenzandoLeerInt)
        {
            // Todavía no cemenzamos a leer el int
            if (c == ' ' || c == '\t')
            {
                // ignoramos
                continue;
            }
            else if (*newLine)
            {
                // nuevo línea, terminamos.
                return true;
            }
            else if (c >= '0' && c <= '9')
            {
                // válido
                buff[idx++] = c;
                comenzandoLeerInt = true;
                *OK = true;
            }
            else if (c == '-')
            {
                // también válido pero el int todavía no es OK
                // porque necesitamos un número después de un -
                buff[idx++] = c;
                comenzandoLeerInt = true;
            }
            else
            {
                // error
                fprintf(stderr, "Found invalid character %c\n", c);
                return false;
            }
        }
        else
        {
            // ya estamos leyendo data
            if (c == ' ' || c == '\t' ||
                c == ',' || *newLine)
            {
                // terminamos
                if (*OK)
                {
                    buff[idx] = '\0';
                    *ll = strtoll(buff, NULL, 10);
                    if (errno != 0)
                    {
                        fprintf(stderr, "Failed to convert %s to long long, error: %s\n", buff, strerror(errno));
                        return false;
                    }
                    return true;
                }
                else
                {
                    // solo podríamos estar aquí si leemos
                    // un '-' y después nada, eso es un error
                    fprintf(stderr, "Found invalid entry \"-\"\n");
                    return false;
                }
            }
            else if (c == '-')
            {
                // un - aquí no es válido porque estámos en el medio
                // de un int.
                fprintf(stderr, "Found \"-\" in the middle of an integer\n");
                return false;
            }
            else if (c >= '0' && c <= '9')
            {
                // válido
                if (idx >= MAX_CHARS)
                {
                    fprintf(stderr, "Integer read was too large to fit into a long long\n");
                    return false;
                }
                buff[idx++] = c;
                *OK = true;
            }
            else
            {
                // error
                fprintf(stderr, "Found invalid character %c\n", c);
                return false;
            }
        }
    }
}

static bool leerComando(FILE *f, bool *eof, Comando *cmd)
{
    // comando más largo es MR, así usamos un buffer de
    // 3 bytes con [2] = NULL
    char buff[3];
    memset(buff, 0, sizeof(buff));

    // read char 1, puede ser M, R o W
    // ignoramos espacio
    while (1)
    {
        int res = fgetc(f);
        if (res == EOF)
        {
            *eof = true;
            return false;
        }

        char c = (char)res;
        if (isspace(c))
        {
            continue;
        }

        // es un carácter no espacio
        // es M, R o W?
        if (c != 'M' &&
            c != 'W' &&
            c != 'R')
        {
            // invalido
            fprintf(stderr, "Read start of invalid command, skipping line.\n");

            // leer hasta '\n'
            skipLine(f);

            return false;
        }

        // es válido
        buff[0] = c;
        break;
    }

    // ahora leer el comando hasta encontramos un espacio
    int idx = 1;
    while (1)
    {
        int res = fgetc(f);
        if (res == EOF)
        {
            *eof = true;
        }
        if (isspace((char)res) ||
            *eof)
        {
            // comando es válido?
            if (strcmp(buff, "R") == 0)
            {
                *cmd = Comando_R;
            }
            else if (strcmp(buff, "W") == 0)
            {
                *cmd = Comando_W;
            }
            else if (strcmp(buff, "MR") == 0)
            {
                *cmd = Comando_MR;
            }
            else
            {
                // invalido
                fprintf(stderr, "Read start of invalid command, skipping line.\n");

                // leer hasta '\n'
                skipLine(f);

                return false;
            }

            return true;
        }

        if (idx > 1)
        {
            fprintf(stderr, "Command too long\n");
            // leer hasta '\n'
            skipLine(f);

            return false;
        }

        buff[idx++] = (char)res;
    }
}

static int leerArchivo(const char *archivo)
{
    FILE *f = fopen(archivo, "r");
    if (f == NULL)
    {
        fprintf(stderr, "%s: No such file or directory\n", archivo);
        return 1;
    }

#ifdef DEBUG
    dumpCache();
#endif


    while (1)
    {
        Comando cmd;
        bool eof = false;
        if (!leerComando(f, &eof, &cmd))
        {
            if (eof)
            {
                break;
            }
        }
        else
        {
            if (cmd == Comando_MR)
            {
                uint32_t mr = get_miss_rate();
                printf("Miss Rate: %u%%\n", mr);
            }
            else
            {
                long long direccion;
                bool OK;
                bool newLine;
                if (!leerLongLong(f, &direccion, &OK, &eof, &newLine))
                {
                    if (eof)
                    {
                        fprintf(stderr, "EOF in middle of command");
                        break;
                    }
                    fprintf(stderr, "Failed to read address\n");
                    if (!newLine)
                    {
                        skipLine(f);
                    }
                    continue;
                }
                else if ((direccion < 0) ||
                         (direccion >= (TAMANO_DE_MEMORIA)))
                {
                    fprintf(stderr, "Address out of bounds\n");
                    if (!newLine)
                    {
                        skipLine(f);
                    }
                    continue;
                }

                if (cmd == Comando_R)
                {
#ifdef DEBUG
                    printf("R %u\n", direccion);
#endif
                    read_byte((int)direccion);
                }
                else
                {
                    if (newLine)
                    {
                        fprintf(stderr, "Comand ended too early");
                        continue;
                    }

                    long long valor;
                    if (!leerLongLong(f, &valor, &OK, &eof, &newLine))
                    {
                        if (eof)
                        {
                            fprintf(stderr, "EOF in middle of command");
                            break;
                        }
                        fprintf(stderr, "Failed to read value\n");
                        if (!newLine)
                        {
                            skipLine(f);
                        }
                        continue;
                    }
                    else if ((valor < 0) ||
                             (valor > 0xFF))
                    {
                        fprintf(stderr, "Value out of bounds\n");
                        if (!newLine)
                        {
                            skipLine(f);
                        }
                        continue;
                    }

#ifdef DEBUG
                    printf("W %u, %u\n", direccion, valor);
#endif
                    write_byte((int)direccion, (uint8_t)valor);
                }
            }
        }
#ifdef DEBUG
        dumpCache();
#endif
    }

#ifdef DEBUG
    dumpCache();
#endif

    fclose(f);

    return 0;
}

int main(int argc, char **argv)
{
    // usamos argv[0] como el nombre del aplicación
    // pero solo queremos el archivo, no la ruta
    const char *nuestroNombre = basename(argv[0]);

    // clear errors
    opterr = 0;

    // parse short options
    while (1)
    {
        // obtener el siguiente argumento
        int option_index = 0;
        int c = getopt_long(argc, argv, "hV", _gLongOptions, &option_index);

        if (c == -1)
        {
            // no hay más
            break;
        }

        switch (c)
        {
            case 'h':
            {
                usage(stdout, nuestroNombre);
                // no seguimos despues de -h
                return 0;
            }
            case 'V':
            {
                printf("%s: Version %u.%u\n", nuestroNombre, MAJOR_VERSION, MINOR_VERSION);
                // no seguimos despues de -V
                return 0;
            }
            case '?':
            {
                if (optopt == 'o')
                {
                    fprintf(stderr, "Option '-%c' requires an argument.\n\n", optopt);
                }
                else if (isprint(optopt))
                {
                    // es un argumento, pero no es uno que esperamos
                    fprintf (stderr, "Unknown option '-%c'.\n\n", optopt);
                }
                else
                {
                    // solo muestra el usage
                }
                usage(stderr, nuestroNombre);
                return 1;
            }
            default:
            {
                usage(stderr, nuestroNombre);
                return 1;
            }
        }
    }

    if (optind == argc)
    {
        fprintf(stderr, "filename is required\n\n");
        usage(stderr, nuestroNombre);
        return 1;
    }

    if ((optind + 1) != argc)
    {
        fprintf(stderr, "Too many arguments\n\n");
        usage(stderr, nuestroNombre);
        return 1;
    }

    init();
    return leerArchivo(argv[optind]);
}
\end{minted}

\end{document}